\section{Anhang}
\label{sec:anhang}

\subsection{Literaturverzeichnis}
\label{Literaturverzeichnis}

Ambrosini, Y. & Merz, L. (2022). Deformationsmessungen Schlipf, Riehen BS [Bachelor-Thesis, Fachhochschule Nordwestschweiz FHNW]. https://www.fhnw.ch/de/studium/architektur-bau-geomatik/bachelor-studiengang-geomatik/bachelor-thesis-2022

Banyai, L. (2011). Rigorous 3D integrated adjustment of GPS baselines, geodetic total station and levelling measurements. In FIG Working Week 2011: Bridging the Gap between Cultures, Marrakech, 18.–22. Mai 2011. International Federation of Surveyors (FIG). https://www.fig.net/resources/proceedings/fig_proceedings/fig2011/papers/ts08f/ts08f_banyai_5363.pdf

França, R. M. de, Klein, I. & Veiga, L. A. K. (2021). The influence of the deflection of the vertical on geodetic surveys in Brazil. Boletim de Ciências Geodésicas, 27(spe), e2021020. https://doi.org/10.1590/s1982-21702021000s00020

GeoConnexion. (2016, 22. März). New Leica LS15 Digital Level provides industry-leading accuracy. https://www.geoconnexion.com/publication-articles/new-leica-ls15-digital-level-provides-industry-leading-accuracy

Leica Geosystems. (o. J.-a). Leica GS18 data sheet [Datenblatt]. https://bynder.hexagon.com/m/47f3f231ddd904d6/original/Leica-GS18-DS-949487-0422-en-LR.pdf

Leica Geosystems. (o. J.-b). Leica Nova TS60 [Datenblatt]. https://leica-geosystems.com/products/total-stations/robotic-total-stations/leica-nova-ts60

Leica Geosystems. (2016). Meeting in the middle: Gotthard Base Tunnel, world’s longest and deepest rail tunnel, benefits from Leica Geosystems measurement technology [Fallstudie]. https://leica-geosystems.com/case-studies/heavy-construction/meeting-in-the-middle

Möser, M. (Hrsg.). (2016). Handbuch Ingenieurgeodäsie: Ingenieurbau (2., neu bearb. und erw. Aufl.). Wichmann.

Niemeier, W. (2008). Ausgleichungsrechnung: Statistische Auswertemethoden (2., überarb. und erw. Aufl.). de Gruyter. https://doi.org/10.1515/9783110206784

Suurmäki, H., Schüler, T., Klügel, T., Kallio, U., Marila, S., Mähler, S., Poutanen, M. & Saari, T. (2022). Datum problem handling in local tie surveys at Wettzell and Metsähovi. In International Association of Geodesy Symposia. Springer. https://doi.org/10.1007/1345_2022_155

swisstopo. (o. J.). LTOP: Präanalyse- und Ausgleichungssoftware für geodätische Netze. Abgerufen am [Datum], von https://www.swisstopo.admin.ch/de/geodaetische-software-ltop

swisstopo. (2004). Geoidmodell der Schweiz CHGeo2004 [Datensatz]. https://www.swisstopo.admin.ch/de/geoid

Vasi, S. & Meyenberg, S. (2014). Deformationsmessung Strain-Zone Le Pont (Jura Decke) [Bachelor-Thesis, Fachhochschule Nordwestschweiz FHNW]. https://www.fhnw.ch/de/studium/architektur-bau-geomatik/bachelor-studiengang-geomatik/bachelor-thesis-2014/media/bth14-01.pdf

Weber, F. (2014). Ceneri-Basistunnel – Portalnetz Zwischenangriff Sigirino [Bachelor-Thesis, Fachhochschule Nordwestschweiz FHNW]. https://www.fhnw.ch/de/studium/architektur-bau-geomatik/bachelor-studiengang-geomatik/bachelor-thesis-2014/media/bth14-04.pdf

Willi, D., Borlat, E. & Deillon, Y. (2024, 15. Mai). Neues Höhensystem der Schweiz: Einleitung, Stand, Roadmap [Präsentation]. GEOSummit 2024, Olten. https://swiss-height-system.heig-vd.ch/wp-content/uploads/2024/05/2024-05-15-GeoSummit.pdf